\section{연습문제}
\label{sec:normal-exercises}

문제는 A, B, C의 세 단계로 나뉜다. 별표 문제는 다음 두 절에 상세 풀이가 있다.

\subsection*{A. 판정과 계산}

\begin{exercise}[*]
$\alpha=\sqrt2+\sqrt3$의 $\Q$-켤레를 모두 구하고 최소다항식을 계산하라. 각 켤레가 $\Q(\sqrt2,\sqrt3)$의 어떤 자동동형으로 얻어지는지 설명하라.
\end{exercise}

\begin{exercise}[*]
$\alpha=\sqrt[3]{2}$라 하자.
\begin{enumerate}[label=(\alph*)]
  \item 모든 $\Q$-매장 $\Q(\alpha)\to\C$를 구하라.
  \item 각 매장의 상을 구하라.
  \item 서로 다른 상들이 생성하는 합성체를 구하라.
\end{enumerate}
\end{exercise}

\begin{exercise}[*]
다음 확장의 정규성을 판정하라.
\[
  \Q(\sqrt[3]{2})/\Q,
  \qquad
  \Q(\sqrt[3]{2},\omega)/\Q,
  \qquad
  \Q(\sqrt2,\sqrt3)/\Q.
\]
정규인 경우 분해체가 되는 다항식을 제시하라.
\end{exercise}

\begin{exercise}
$K$의 표수가 $2$가 아니고 $a\in K$가 제곱이 아니라고 하자. $K(\sqrt a)/K$가 정규임을 최소다항식과 체매장의 두 관점에서 각각 증명하라.
\end{exercise}

\begin{exercise}[*]
\[
  \alpha=\sqrt{2+\sqrt2},
  \qquad
  L=\Q(\alpha)
\]
라 하자.
\begin{enumerate}[label=(\alph*)]
  \item $\alpha$의 최소다항식을 구하라.
  \item 모든 $\Q$-켤레가 $L$에 속함을 보여라.
  \item $L/\Q$의 자동동형 네 개를 $\alpha$의 상으로 기술하라.
\end{enumerate}
\end{exercise}

\begin{exercise}
$\alpha=\sqrt[4]{2}$라 하자. $\Q(\alpha)/\Q$가 정규가 아님을 보이고 정상폐포와 그 차수를 구하라.
\end{exercise}

\begin{exercise}[*]
$\alpha=\sqrt[5]{2}$이고 $\zeta_5$가 원시 오차단위근이라 하자. $\Q(\alpha)/\Q$의 정상폐포가
\[
  \Q(\alpha,\zeta_5)
\]
임을 보이고 그 차수를 구하라.
\end{exercise}

\begin{exercise}
유한확장 $\F_{q^n}/\F_q$가 정규임을 다음 두 방법으로 증명하라.
\begin{enumerate}[label=(\alph*)]
  \item $X^{q^n}-X$의 분해체를 이용한다.
  \item 상대 프로베니우스와 모든 $\F_q$-매장을 이용한다.
\end{enumerate}
\end{exercise}

\begin{exercise}[*]
$L=\Q(\sqrt2)$에서 $f=X^4-2$를 인수분해하라. $L/\Q$의 비자명한 자동동형이 두 기약인수를 서로 바꿈을 확인하라.
\end{exercise}

\begin{exercise}
$K=\F_p(t^p)$, $L=\F_p(t)$라 하자. $L/K$가 정규이고 $\Aut_K(L)$가 자명함을 보인 뒤
\[
  L^{\Aut_K(L)}
\]
를 구하라.
\end{exercise}

\subsection*{B. 핵심 정리의 증명}

\begin{exercise}[*]
$K$의 대수적 폐포 $\overline K$ 안의 대수적 원소 $\alpha,\beta$에 대해 다음 조건들의 동치를 증명하라.
\begin{enumerate}[label=(\roman*)]
  \item $m_{\alpha,K}=m_{\beta,K}$.
  \item $K(\alpha)\cong_K K(\beta)$이며 $\alpha$를 $\beta$로 보낼 수 있다.
  \item 어떤 $\sigma\in\Aut_K(\overline K)$가 $\sigma(\alpha)=\beta$를 만족한다.
\end{enumerate}
\end{exercise}

\begin{exercise}
대수적 확장 $L/K$에 대해 다음 조건들의 동치를 직접 증명하라.
\begin{enumerate}[label=(\roman*)]
  \item 모든 $K$-매장 $L\to\Omega$의 상이 $L$이다.
  \item $L$은 $K[X]$의 다항식족의 분해체이다.
  \item 기약다항식이 $L$ 안에 한 근을 가지면 $L$에서 완전히 분해된다.
\end{enumerate}
\end{exercise}

\begin{exercise}[*]
확장탑 $K\subseteq E\subseteq M$에서 $M/K$가 정규이면 $M/E$가 정규임을 증명하라. 이어서
\[
  \Q\subseteq\Q(\sqrt2)\subseteq\Q(\sqrt[4]{2})
\]
를 사용하여 정규성이 추이적이지 않음을 보여라.
\end{exercise}

\begin{exercise}
$L/K$가 정규이고 $E/K$가 같은 대수적 폐포 안의 대수적 확장이라 하자. $LE/E$가 정규임을 분해체의 관점에서 증명하라.
\end{exercise}

\begin{exercise}[*]
하나의 대수적 폐포 안에서 $L_1/K$와 $L_2/K$가 정규라고 하자. 다음 확장들이 정규임을 증명하라.
\[
  L_1L_2/K,
  \qquad
  (L_1\cap L_2)/K.
\]
\end{exercise}

\begin{exercise}
$L=K(\alpha_i:i\in I)$가 대수적 확장이라 하자. 각 $\alpha_i$의 최소다항식들의 분해체가 $L/K$의 정상폐포임을 증명하라.
\end{exercise}

\begin{exercise}
대수적 폐포 $\Omega$ 안에서
\[
  N=K\bigl(\sigma(L):\sigma\in\Hom_K(L,\Omega)\bigr)
\]
가 $L/K$의 정상폐포임을 증명하라.
\end{exercise}

\begin{exercise}[*]
$L/K$가 유한 정규확장임을 가정하라.
\begin{enumerate}[label=(\alph*)]
  \item 모든 $K$-매장 $L\to\Omega$가 $L$의 자동동형임을 보여라.
  \item $|\Aut_K(L)|=[L:K]_{\mathrm s}$임을 보여라.
  \item $L/K$가 분리일 필요충분조건이 $|\Aut_K(L)|=[L:K]$임을 보여라.
\end{enumerate}
\end{exercise}

\begin{exercise}
$L/K$가 정규이고 $f\in K[X]$가 일계수 기약이라고 하자. $f$의 $L[X]$에서의 일계수 기약인수들이 $\Aut_K(L)$ 아래에서 하나의 궤도를 이루며 모두 같은 차수임을 증명하라.
\end{exercise}

\begin{exercise}
$L/K$가 정규 분리확장이라고 하자. 모든 $K$-자동동형이 고정하는 원소는 $K$에 속함을, 켤레집합과 분리성을 사용하여 증명하라.
\end{exercise}

\subsection*{C. 구조와 응용}

\begin{exercise}[*]
$\alpha=\sqrt[4]{2}$, $N=\Q(\alpha,i)$라 하자.
\begin{enumerate}[label=(\alph*)]
  \item 모든 $\Q$-자동동형을 $\alpha$와 $i$의 상으로 기술하라.
  \item 자동동형이 정확히 여덟 개임을 보여라.
  \item $\alpha$의 궤도와 안정자를 구하라.
\end{enumerate}
\end{exercise}

\begin{exercise}
$N=\Q(\sqrt[3]{2},\omega)$의 모든 $\Q$-자동동형을 구하고 $\sqrt[3]{2}$와 $\omega$의 궤도를 각각 계산하라.
\end{exercise}

\begin{exercise}[*]
$L=\Q(\sqrt2,\sqrt3)$, $G=\Aut_\Q(L)$라 하자.
\begin{enumerate}[label=(\alph*)]
  \item $G$의 원소를 두 제곱근의 부호 변화로 나타내라.
  \item $\sqrt2+\sqrt3$의 안정자를 구하라.
  \item $\sqrt2$의 안정자와 궤도를 구하고 궤도-안정자 공식을 확인하라.
\end{enumerate}
\end{exercise}

\begin{exercise}
문제 \thechapter.5의 체 $L$에서 $\alpha\mapsto\sqrt{2-\sqrt2}$인 자동동형의 위수가 $4$임을 보이고 $\Aut_\Q(L)$의 군 구조를 결정하라.
\end{exercise}

\begin{exercise}[*]
$p$가 홀수이고 $K=\F_p(s,t)$라 하자. $\alpha^p=s$, $\beta^2=t$인 원소들을 첨가하여
\[
  L=K(\alpha,\beta)
\]
라 하자.
\begin{enumerate}[label=(\alph*)]
  \item $L/K$가 정규임을 보여라.
  \item $\Aut_K(L)$을 구하라.
  \item 고정체 $K_{\mathrm i}=L^{\Aut_K(L)}$와 분리폐포 $K_{\mathrm s}$를 구하라.
  \item $L=K_{\mathrm s}K_{\mathrm i}$를 확인하라.
\end{enumerate}
\end{exercise}

\begin{exercise}
순수비분리 대수적 확장 $L/K$의 모든 $K$-자동동형이 항등사상임을 증명하라. 역이 일반적으로 성립하는지 논하라.
\end{exercise}

\begin{exercise}[*]
$L=\Q(\sqrt2,i)$ 위에서 $f=X^4-2$를 기약인수들의 곱으로 분해하라. 각 인수가 기약임을 보이고 $\Aut_\Q(L)$의 어떤 원소가 두 인수를 서로 바꾸는지 찾아라.
\end{exercise}

\begin{exercise}
$L/K$가 유한 정규확장이고 $\alpha\in L$라 하자. 안정자
\[
  G_\alpha=\{\sigma\in\Aut_K(L):\sigma(\alpha)=\alpha\}
\]
에 대해
\[
  [K(\alpha):K]_{\mathrm s}
  =\frac{|\Aut_K(L)|}{|G_\alpha|}
\]
를 증명하고, $\alpha$가 순수비분리일 때 식이 어떻게 되는지 설명하라.
\end{exercise}

\begin{exercise}[*]
$L/K$가 유한 정규확장이고 $K_{\mathrm s}$와 $K_{\mathrm i}$가 각각 분리폐포와 자동동형군의 고정체라 하자. 다음 순서로
\[
  L=K_{\mathrm s}K_{\mathrm i}
\]
를 증명하라.
\begin{enumerate}[label=(\alph*)]
  \item $K_{\mathrm s}\cap K_{\mathrm i}=K$임을 보여라.
  \item $[L:K_{\mathrm i}]=[K_{\mathrm s}:K]$임을 보여라.
  \item 합성체의 차수를 비교하여 결론을 얻어라.
\end{enumerate}
\end{exercise}

\begin{exercise}
유한 단순확장 $L=K(\alpha)$에 대해 다음 조건들이 동치임을 증명하라.
\begin{enumerate}[label=(\roman*)]
  \item $L/K$는 정규이다.
  \item $\alpha$의 모든 $K$-켤레가 $L$에 들어 있다.
  \item 모든 $K$-켤레 $\beta$에 대해 $K(\beta)=L$이다.
\end{enumerate}
\end{exercise}
